\documentclass[11pt]{article}
\usepackage[margin=1in]{geometry}
\usepackage{amsmath,amssymb,amsthm}
\usepackage[numbers,sort&compress]{natbib}
\usepackage{hyperref}

\newtheorem{theorem}{Theorem}
\newtheorem{proposition}[theorem]{Proposition}
\newtheorem{lemma}[theorem]{Lemma}
\newtheorem{corollary}[theorem]{Corollary}
\theoremstyle{definition}
\newtheorem{definition}[theorem]{Definition}
\newtheorem{remark}[theorem]{Remark}

\newcommand{\Is}{I^-_{\mathrm s}}
\newcommand{\Isp}{I^+_{\mathrm s}}
\newcommand{\Det}{\mathrm{Det}}

\title{Unreal Is Not Open:\\ prediction, horizons and the future in
relativistic becoming}
\author{Andrew Impellitteri}
\date{25 October 2026 (draft 2)}

\begin{document}
\maketitle

\begin{abstract}
Local theories of relativistic becoming make what is real from an event its
causal past. Does the unreality of the future from an event make the future
open from it? We separate the two questions. Call an event determined from
$b$ if every past-inextendible causal curve from it meets $b$'s
chronological past; under deterministic laws with the domain-of-dependence
property, such events are fixed by what is real from $b$. The notion is
essentially Geroch's domain of prediction. We characterize it by a condition
on the slices of any Cauchy time function, show that it is conformally
invariant, and compute it for the standard cosmological models. In
Minkowski spacetime, spatially flat and open big-bang universes, the closed
radiation universe and de Sitter spacetime, nothing outside an event's past
is ever determined from it. In the closed dust model, nothing is while the
universe expands and everything is from every event after maximum expansion,
the known moment at which the particle horizon disappears. So unreality and
openness come apart in both directions, and, given the laws, which
combination holds is a matter of global causal structure and can change with
cosmic epoch. A closing remark records what the corresponding standpoint
logic does and does not show about cosmological horizons.
\end{abstract}

\section{Introduction}

Stein's theorem says that in Minkowski spacetime a transitive,
symmetry-invariant relation of becoming that makes some past event definite
is the causal order or trivial \citep{Stein1968,Stein1991,CH1995}. A
companion paper \citep{SteinGR} extends this to general relativity under a
locality principle: what has become for an event $b$ always includes its
chronological past and never goes beyond its causal past. On views that read
``has become for $b$'' as ``is real from $b$'' \citep{Fine2005,Lipman2023},
the future of $b$ is not real from $b$.

It is tempting to infer that the future is \emph{open} from $b$. The
inference fails, and the way it fails is instructive. Openness in the sense
relevant here is a matter of what the real fixes, and that depends on the
laws and on the global causal structure. The point is old in outline: a
growing block \citep{Broad1923} combined with deterministic laws makes the
future unreal yet fixed, and Rak\'ic's sea-fight case turns on an event
outside an observer's past that a decision inside it appears to settle
\citep{Rakic1997}. What is missing is a precise account, for relativistic
spacetimes, of when what is real from an event fixes anything beyond itself.
Geroch's domain of prediction supplies the tool \citep{Geroch1977}, and
Manchak's work on prediction its limits \citep{Manchak2008}. We use them to
answer three questions:
\begin{enumerate}
\item Exactly which events are fixed, under deterministic laws, by what is
real from $b$ (Theorem~\ref{thm:char})?
\item In the standard cosmological models, how much is (Section~\ref{sec:examples})?
\item What follows for the relation between unreality and openness
(Section~\ref{sec:phil})?
\end{enumerate}

\paragraph{Conventions.} Spacetimes are smooth, connected, time-oriented and
globally hyperbolic. $\Is(p)$ is the strict chronological past (via timelike
curves), $J^-(p)$ the causal past; $x\le y$ means $x\in J^-(y)$, and $x\ll y$
means $x\in\Is(y)$. We use that $\le$ is closed, that $J^-(x)\cap S$ is
compact for every Cauchy surface $S$, push-up ($y\le z\ll b$ implies $y\ll b$),
and that $M=\Is(S)\cup S\cup\Isp(S)$ for a Cauchy surface $S$ \citep{MS2008}.
A \emph{Cauchy time function} is a continuous surjection $t:M\to(\alpha,\beta)$,
$-\infty\le\alpha<\beta\le\infty$, strictly increasing along future-directed
causal curves, whose level sets $S_c=t^{-1}(c)$ are Cauchy surfaces
\citep{BS2005}.

\section{Determined from an event}\label{sec:det}

\begin{definition}
An event $x$ is \emph{determined from} $b$ if every past-inextendible causal
curve from $x$ meets $\Is(b)$. Write $\Det(b)$ for the set of such events.
Clearly $\Is(b)\subseteq\Det(b)$.
\end{definition}

\begin{remark}[What the definition does and does not say]\label{rem:laws}
$\Det(b)$ is a property of the spacetime and the event. It becomes a claim
about what is fixed only given laws with the \emph{domain-of-dependence
property}: data on a region fix the fields at every event all of whose
past-inextendible causal curves meet the region. Under such laws we call the
events of $\Det(b)$ \emph{nomologically settled from $b$}. This is not
Belnap's notion of settledness, on which a proposition is settled at a point
event if it holds in every possible history through it \citep{Belnap1992};
under indeterministic laws the two come apart, and $\Det(b)$ says nothing
about them. The domain-of-dependence property is also a locality assumption
on the dynamics. We use $\Is(b)$, which every local becoming theory counts as
real from $b$, rather than $J^-(b)$; the characterization below uses that
$\Is(b)$ is open.
\end{remark}

\paragraph{Relation to Geroch.} Geroch's domain of prediction $P(q)$ is the set
of $x$ such that every past-directed timelike curve from $x$ without past
endpoint enters $I^-(q)$, but $I^-(x)\not\subseteq I^-(q)$ \citep[sec.~3]{Geroch1977}.
$\Det(b)$ is the causal-curve analogue of the union of $P(b)$ with the
events whose pasts lie in that of $b$; the two differ at most on null
boundaries. Geroch gives an equivalent form by achronal surfaces in $I^-(q)$
\citep[sec.~3]{Geroch1977}; Theorem~\ref{thm:char} is a slicing form for
globally hyperbolic spacetimes.

\begin{theorem}\label{thm:char}
Let $t:M\to(\alpha,\beta)$ be a Cauchy time function. Then $x\in\Det(b)$ if
and only if there is $c\in(\alpha,t(x))$ with $J^-(x)\cap S_c\subseteq\Is(b)$.
\end{theorem}

\begin{proof}
($\Leftarrow$) A past-inextendible causal curve from $x$ crosses the Cauchy
surface $S_c$, at a point of $J^-(x)\cap S_c\subseteq\Is(b)$.

($\Rightarrow$) Suppose that for every $c\in(\alpha,t(x))$ the set
$K_c=(J^-(x)\cap S_c)\setminus\Is(b)$ is nonempty; it is compact, since
$\Is(b)$ is open. Fix $t(x)>c_1>c_2>\dots$ with $c_n\to\alpha$, and let
\[P_n=\{(z_1,\dots,z_n)\in K_{c_1}\times\dots\times K_{c_n}:
x\ge z_1\ge z_2\ge\dots\ge z_n\}.\]
$P_n$ is compact, since $\le$ is closed. It is nonempty: pick $y\in K_{c_n}$
and a causal curve from $x$ to $y$; it crosses each $S_{c_i}$, $i<n$, at some
$z_i\le x$, and $z_i\notin\Is(b)$, for otherwise $y\le z_i\ll b$ would give
$y\ll b$. The maps $P_{n+1}\to P_n$ forgetting the last coordinate are
continuous, so the inverse limit of the nonempty compact spaces $P_n$ is
nonempty, giving $x\ge z_1\ge z_2\ge\cdots$ with $z_n\in K_{c_n}$. Join
consecutive terms by causal curves and concatenate. Each segment avoids
$\Is(b)$ by push-up, as before. Along the concatenation $t\to\alpha$, so it has
no past endpoint in $M$: it is a past-inextendible causal curve from $x$
missing $\Is(b)$.
\end{proof}

\begin{corollary}\label{cor:cauchy}
If $\Is(b)$ contains a Cauchy surface $S$, then $\Det(b)=M$.
\end{corollary}

\begin{proof}
If $x\in\Is(S)\cup S$, then $x\le s\ll b$ for some $s\in S$, so $x\in\Is(b)$.
If $x\in\Isp(S)$, extend a past-inextendible causal curve from $x$ to an
inextendible one; it meets $S$, and not to the future of $x$, since then
$s'\ll x\le s$ for some $s,s'\in S$, contradicting achronality. So the curve
from $x$ meets $S\subseteq\Is(b)$.
\end{proof}

Since $\Det(b)$ is defined by causal curves alone, it depends only on the
conformal class of the metric; conformal diagrams decide it.

\section{The standard models}\label{sec:examples}

\begin{lemma}\label{lem:balls}
Let $(\Sigma,d)$ be $\mathbb R^3$, hyperbolic space $H^3$, or the unit sphere
$S^3$ with $r_2\le\pi$. Then the closed ball $\bar B(X,r_1)$ lies in the open
ball $B(Y,r_2)$ if and only if $d(X,Y)+r_1<r_2$.
\end{lemma}

\begin{proof}
($\Leftarrow$) is the triangle inequality. ($\Rightarrow$) Suppose
$d(X,Y)+r_1\ge r_2$. If $d(X,Y)+r_1\le\pi$ (always so on $\mathbb R^3$ and
$H^3$), follow a minimizing geodesic from $Y$ through $X$ (any geodesic from
$X$ if $X=Y$) a further distance $r_1$; the endpoint lies in $\bar B(X,r_1)$ at
distance $d(X,Y)+r_1\ge r_2$ from $Y$. If $d(X,Y)+r_1>\pi$ on $S^3$, the
antipode of $Y$ is at distance $\pi-d(X,Y)<r_1$ from $X$ and $\pi\ge r_2$ from
$Y$.
\end{proof}

\begin{proposition}\label{prop:conformal}
Let $M$ be conformal to $(\eta_-,\eta_+)\times\Sigma$ with metric
$-d\eta^2+d\Sigma^2$, where $\Sigma$ is $\mathbb R^3$, $H^3$, or the unit $S^3$
with $\eta_+-\eta_-\le\pi$. Then $\Det(b)=\Is(b)$ for every $b$.
\end{proposition}

\begin{proof}
Use $t=\eta$. For $x=(\eta_x,X)$, $b=(\eta_b,Y)$: $x\ll b$ iff
$d(X,Y)<\eta_b-\eta_x$ (on $S^3$, $\eta_b-\eta_x<\pi$). For $\eta_-<c<\eta_x$,
$J^-(x)\cap S_c=\bar B(X,\eta_x-c)$ and $\Is(b)\cap S_c=B(Y,\eta_b-c)$, with
$\eta_b-c<\pi$ on $S^3$. By Lemma~\ref{lem:balls} the inclusion holds iff
$d(X,Y)<\eta_b-\eta_x$, i.e.\ iff $x\ll b$. Apply Theorem~\ref{thm:char}.
\end{proof}

This covers Minkowski spacetime, spatially flat and open FLRW universes with
a big bang (conformal to $\eta>0$), the closed radiation universe
($a\propto\sin\eta$, $\eta\in(0,\pi)$) and global de Sitter spacetime ($\Sigma=S^3$,
$\eta\in(-\pi/2,\pi/2)$) \citep{HE1973}. In none of them is anything outside an
event's chronological past determined from it.

\begin{proposition}\label{prop:dust}
Let $M$ be conformal to $(\eta_-,\eta_+)\times S^3$ with $\eta_+-\eta_->\pi$.
Then $\Det(b)=\Is(b)$ if $\eta_b\le\eta_-+\pi$, and $\Det(b)=M$ if
$\eta_b>\eta_-+\pi$. In the Einstein static universe ($\eta_\pm=\pm\infty$),
$\Det(b)=M$ for every $b$.
\end{proposition}

\begin{proof}
If $\eta_b>\eta_-+\pi$, take $c\in(\eta_-,\eta_b-\pi)$: $\Is(b)\cap S_c$ is an open
ball of radius greater than $\pi$, hence all of $S_c$, and
Corollary~\ref{cor:cauchy} applies. If $\eta_b\le\eta_-+\pi$, every $c>\eta_-$
has $\eta_b-c<\pi$, and the proof of Proposition~\ref{prop:conformal} applies.
\end{proof}

For the closed dust FLRW model, $a(\eta)=\tfrac12a_m(1-\cos\eta)$ with
$\eta\in(0,2\pi)$, the threshold is $\eta=\pi$, the moment of maximum
expansion. That the particle horizon of the closed dust model disappears at
maximum expansion, so that an observer can then in principle see the whole
universe, is known \citep{GronJohannesen2011,BolotinTanatarov2013}; the
Einstein static case is Geroch's \citep[Fig.~5]{Geroch1977}. What
Proposition~\ref{prop:dust} adds is the consequence for determination: with
Corollary~\ref{cor:cauchy}, the vanishing of the particle horizon is exactly
the condition under which what is real from $b$ fixes everything. The result
fits Manchak's theorem that prediction of events to the future of an observer
requires a compact Cauchy surface \citep{Manchak2008}: only spatially closed
models can cross the threshold, and not all of them do.

\section{Unreality and openness}\label{sec:phil}

Put together with the companion theorem, these results separate two
properties of an event $x$ relative to a standpoint $b$:
\begin{itemize}
\item \emph{real from $b$}: $x$ is in $b$'s fragment, which a local becoming
theory fixes to lie between $\Is(b)$ and $J^-(b)$;
\item \emph{settled from $b$}: under laws with the domain-of-dependence
property, $x\in\Det(b)$.
\end{itemize}
Every real event (with the possible exception of the null boundary) is
settled. The converse fails in two ways. In the flat, open, radiation and de
Sitter models nothing beyond the fragment is settled; in the closed dust model
after maximum expansion, and in the Einstein static universe, everything is.
So:

\begin{enumerate}
\item \emph{Unreality does not entail openness.} An event can be unreal from $b$
and settled from $b$: in the closed dust model, every future event is so from
every event after maximum expansion. This is the relativistic form of the
growing block combined with determinism, and on our reading it is also where
Rak\'ic's intuition belongs \citep{Rakic1997}: her sea-fight is, at most,
settled from the observer, not real from her. In Minkowski spacetime it is not
even settled: by Proposition~\ref{prop:conformal}, nothing outside the
observer's past is, whatever lies in that past.
\item \emph{Determinism does not entail settledness from each event.} Under
deterministic laws the whole spacetime is fixed by data on any Cauchy surface,
but no single event's past contains a Cauchy surface in the flat, open,
radiation or de Sitter models. From each standpoint, the future is then
unsettled as well as unreal.
\item \emph{Which combination holds is a matter of global causal structure and
can change with cosmic epoch.} Given the laws, the answer to ``is the future
open from here?'' can differ between two events of the same spacetime.
\end{enumerate}

Three limits. First, this openness is relational and informational: an event
outside $\Det(b)$ is still fixed by its own past, and a determinist eternalist
can accept everything in this section. What the section shows is that the
standpoint theorist's claim that the future is unreal from $b$ neither
secures nor excludes its being fixed from $b$; an ontically open future needs
indeterministic laws, about which $\Det(b)$ is silent. Second, the closed
dust result concerns an idealized model with deterministic, classical
dynamics; it is not a claim about our universe, which on current evidence is
not of that kind. Third, Manchak shows that an observer who does not already
know the spacetime cannot verify such predictions \citep{Manchak2008};
$\Det(b)$ is a nomological notion given $M$, not an epistemic one.

\section{A remark on standpoint logic}\label{sec:logic}

Read $\Box A$ at $b$ as ``at every event $x\le b$, $A$'' on the frame $(M,R)$
with $bRx$ iff $x\le b$. If $A$ is a tenseless claim about what occurs at $x$,
this is the reflexive past modality of Diodorean tense logic, and the
standpoint reading of Fine and Lipman \citep{Fine2005,Lipman2023} adds no
formal content. $R$ is reflexive and transitive, so S4 is valid, and in
Minkowski spacetime the valid formulas are exactly those of S4.2
\citep{Goldblatt1980}. On reflexive transitive frames the axiom .2,
$\Diamond\Box A\to\Box\Diamond A$, is valid iff the frame is directed
\citep{BdRV2001}; on the subframe generated by $b$, directedness says that any
two events in $J^-(b)$ have a common causal past.

In a model conformal to $(\eta_-,\eta_+)\times\mathbb R^3$, two events have a
common past iff $|X-Y|<\eta_x+\eta_y-2\eta_-$, so if $\eta_->-\infty$ the
directedness fails in the past of every event and .2 is not valid on any
generated subframe. This is the particle-horizon structure of Rindler
\citep{Rindler1956} in modal form, and it holds equally with or without an
inflationary era of finite conformal duration. The horizon problem of
cosmology \citep{Guth1981} is therefore not the failure of .2, which holds on
both sides of it, but the failure of particular instances: pairs of points
on our last-scattering surface, at comoving separation up to
$2(\eta_0-\eta_{\rm ls})$, have no common past unless
$\eta_0-\eta_{\rm ls}<\eta_{\rm ls}-\eta_-$, which fails without inflation and
holds with enough of it. Two cautions. Under the view's persistence principle,
on which what holds from an event holds from every later one, valuations are
constrained along the future direction, and the frame property that matters
may concern common futures, which these models have when $\eta_+=\infty$; the logic of persistent
valuations is a separate question. And Lewis and van der Schaaf study causal
modal logics of general spacetimes \citep{vdSchaaf2026}; we have not checked
how their results bear on the past-directed frames used here. Whether S4 is
exactly the logic of all past-directed globally hyperbolic frames we leave
open.

\section{Conclusion}

The unreality of the future from an event, which local theories of
relativistic becoming deliver, neither implies nor excludes its being fixed
from that event. Which holds is decided, given deterministic local laws, by a
slicing condition on Geroch's domain of prediction (Theorem~\ref{thm:char}),
and it is computable from conformal diagrams. In most standard models nothing
beyond an event's past is fixed from it; in the closed dust model everything
is, from every event after maximum expansion. The question whether the
future is open from a standpoint thus divides into a question about the laws
and a question about global causal structure, and only the second is settled
by relativity.

\begin{thebibliography}{99}
\bibitem{BdRV2001} P. Blackburn, M. de Rijke and Y. Venema, \emph{Modal Logic}
(Cambridge University Press, 2001).
\bibitem{Belnap1992} N. Belnap, Branching space-time, \emph{Synthese} 92 (1992)
385--434, doi:10.1007/BF00414289.
\bibitem{BolotinTanatarov2013} Yu. L. Bolotin and I. V. Tanatarov, A thousand
problems in cosmology: horizons, arXiv:1310.6329 (2013).
\bibitem{Broad1923} C. D. Broad, \emph{Scientific Thought} (Kegan Paul, 1923).
\bibitem{BS2005} A. N. Bernal and M. S\'anchez, Smoothness of time functions
and the metric splitting of globally hyperbolic spacetimes, \emph{Commun.\
Math.\ Phys.} 257 (2005) 43--50.
\bibitem{CH1995} R. Clifton and M. Hogarth, The definability of objective
becoming in Minkowski spacetime, \emph{Synthese} 103 (1995) 355--387.
\bibitem{Fine2005} K. Fine, Tense and reality, in \emph{Modality and Tense}
(Oxford University Press, 2005) 261--320.
\bibitem{Geroch1977} R. Geroch, Prediction in general relativity, in J. Earman,
C. Glymour and J. Stachel (eds.), \emph{Foundations of Space-Time Theories},
Minnesota Studies in the Philosophy of Science 8 (University of Minnesota
Press, 1977) 81--93.
\bibitem{Goldblatt1980} R. Goldblatt, Diodorean modality in Minkowski
spacetime, \emph{Studia Logica} 39 (1980) 219--236, doi:10.1007/BF00370321.
\bibitem{GronJohannesen2011} \O. Gr\o n and S. Johannesen, FRW universe models in
conformally flat spacetime coordinates. III: universe models with positive
spatial curvature, \emph{Eur.\ Phys.\ J. Plus} 126 (2011) 30,
arXiv:1103.4753.
\bibitem{Guth1981} A. H. Guth, Inflationary universe: a possible solution to
the horizon and flatness problems, \emph{Phys.\ Rev.\ D} 23 (1981) 347--356.
\bibitem{HE1973} S. W. Hawking and G. F. R. Ellis, \emph{The Large Scale
Structure of Space-Time} (Cambridge University Press, 1973).
\bibitem{Lipman2023} M. A. Lipman, Standpoints: a study of a metaphysical
picture, \emph{J.\ Phil.} 120 (2023) 117--138.
\bibitem{Manchak2008} J. B. Manchak, Is prediction possible in general
relativity?, \emph{Found.\ Phys.} 38 (2008) 317--321,
doi:10.1007/s10701-008-9204-6.
\bibitem{MS2008} E. Minguzzi and M. S\'anchez, The causal hierarchy of
spacetimes, in \emph{Recent Developments in Pseudo-Riemannian Geometry}
(EMS, 2008) 299--358, arXiv:gr-qc/0609119.
\bibitem{Rakic1997} N. Raki\'c, Past, present, future, and special relativity,
\emph{Brit.\ J. Phil.\ Sci.} 48 (1997) 257--280.
\bibitem{Rindler1956} W. Rindler, Visual horizons in world-models,
\emph{Mon.\ Not.\ R. Astron.\ Soc.} 116 (1956) 662--677.
\bibitem{Stein1968} H. Stein, On Einstein--Minkowski space--time, \emph{J.\
Phil.} 65 (1968) 5--23.
\bibitem{Stein1991} H. Stein, On relativity theory and openness of the future,
\emph{Phil.\ Sci.} 58 (1991) 147--167.
\bibitem{SteinGR} A. Impellitteri, Stein's theorem in general relativity:
locality, not symmetry, fixes what has become, preprint (2026).
\bibitem{vdSchaaf2026} M. Lewis and N. van der Schaaf, Some results on causal
modalities in general spacetimes, arXiv:2601.14029v2 (2026).
\end{thebibliography}

\end{document}
